\documentclass[11pt]{article}

\usepackage[margin=0.9in]{geometry}
\usepackage[T1]{fontenc}
\usepackage[utf8]{inputenc}
\usepackage{lmodern}
\usepackage{array}
\usepackage{booktabs}
\usepackage{hyperref}
\usepackage{parskip}
\usepackage{ragged2e}
\usepackage{setspace}
\usepackage[table]{xcolor}
\usepackage[natbibapa]{apacite}

\definecolor{PrimaryRed}{HTML}{8B1E1E}
\definecolor{SoftBlue}{HTML}{DDEEFF}
\definecolor{SoftGray}{HTML}{F4F4F4}

\hypersetup{colorlinks=true,linkcolor=black,urlcolor=black,citecolor=black}
\setlength{\parindent}{0pt}
\setstretch{1.08}
\bibliographystyle{apacite}
\newcommand{\reviewsection}[1]{\vspace{0.07in}\noindent{\normalsize\bfseries\color{PrimaryRed}#1}\par\vspace{0.02in}}

\begin{document}

\hypersetup{pageanchor=false}
\begin{titlepage}
\newgeometry{margin=0.7in}
\thispagestyle{empty}

\begin{center}
{\large University of Rochester\\}
{\normalsize Warner School of Education\\[0.15in]}

{\color{PrimaryRed}\rule{\textwidth}{1.4pt}}\\[0.12in]
{\Huge\bfseries Who Gets to Name,\\[0.07in]
Interpret, and Design?}\\[0.15in]
{\huge Understanding Developmental Difference\\[0.07in]
Through Labels, Lived Experience,\\[0.07in]
and Educational Environments}\\[0.18in]
{\Large EDE 448 --- Module 1 Journal}\\[0.12in]

\vfill

{\color{PrimaryRed}\rule{0.78\textwidth}{0.7pt}}\\[0.22in]

\colorbox{SoftBlue}{%
\parbox{0.88\textwidth}{%
\centering\small\vspace{0.02in}
\textbf{From Deficit Interpretation to Access Design}\\[0.05in]
A reflective analysis of how diagnosis, neurodiversity, media, schooling, family knowledge, lived experience, and classroom design shape what the world believes about developmental difference.
}%
}
\end{center}

\vfill

\begin{center}
\renewcommand{\arraystretch}{1.18}
\setlength{\tabcolsep}{4pt}
\begin{tabular}{>{\bfseries\RaggedLeft}p{1.65in} p{5.0in}}
Student & Piter Zacari Garcia Bautista \\
Assignment & Module 1 Journal \\
Course & EDE 448 --- Communication and Positive Behavioral Supports for Students with Autism and Other Complex Needs \\
Essential Question & How does the world come to understand developmental differences? \\
Sources & Cole and Telesford; Ortega; National Education Association; \textit{Autism: Insight from Inside} \\
Interpretive Lenses & Puzzle Plan and EQUITAS \\
Submission Date & July 2026 \\
\end{tabular}
\end{center}

\vfill

\begin{center}
\begin{minipage}{0.92\textwidth}
\centering\itshape
A label can open access, but it becomes harmful when adults allow it to replace curiosity.\\
{\color{PrimaryRed}\rule{4in}{0.8pt}}
\end{minipage}
\end{center}

\vfill

\begin{center}
\begin{minipage}{0.95\textwidth}
\small
This journal argues that developmental difference is simultaneously neurological, lived, relational, cultural, and socially interpreted. Rather than separating course ideas from experience, each section uses my lived experience and K--5 technology teaching placement to interpret the readings and media. The analysis asks how richer evidence can move educators from labeling behavior toward designing environments that reveal competence, protect agency, and support belonging.
\end{minipage}
\end{center}

\restoregeometry
\end{titlepage}
\hypersetup{pageanchor=true}
\pagenumbering{arabic}

\reviewsection{What I Noticed Before I Had the Language}

Before these readings, I already knew what it felt like to be interpreted through fragments. As a neurodivergent person, I have lived through institutions that noticed missed deadlines, uneven performance, communication differences, or difficulty navigating an inaccessible process before they noticed effort, context, intelligence, or persistence. That history shapes how I now look at students. When a child hesitates, repeats a question, moves away from a group, keeps a familiar device close, or needs more time, I do not see a complete explanation. I see a clue.

That is the most important shift Module 1 helped me name. The world does not encounter developmental difference directly. It encounters behavior, language, stories, reports, family knowledge, diagnostic criteria, media images, and school routines. Medicine and psychology name patterns. Schools translate those names into eligibility, instruction, discipline, and support. Families contribute histories that may never appear in an evaluation. Autistic people and advocacy communities challenge the meanings attached to those labels. Each system reveals something, but each can also hide something.

My central learning is therefore personal as well as academic: developmental difference is neurological, lived, relational, cultural, and socially interpreted. No single system gets to explain the whole person.

\reviewsection{When a Label Opens a Door---and When It Starts Closing One}

In my elementary technology teaching placement, I repeatedly saw students whose competence was easy to miss because of what surrounded the task. A learner might understand the coding concept but struggle with a long verbal direction, an unfamiliar login path, reading load, noise, peer control of the device, or uncertainty about what finished work should look like. When I modeled the exact tool path, used a visible first/next/last sequence, broke the task into smaller steps, or gave additional processing time, participation changed. The academic idea had not suddenly entered the student's mind. The environment had stopped hiding what was already there.

That experience is what made \citet{cole2023autism} meaningful to me. Their clinical and educational language can be useful because it gives families and professionals a shared vocabulary for social communication, sensory processing, routines, focused interests, and support needs. A diagnosis can open access to services and help adults recognize patterns that might otherwise be dismissed. At the same time, the reading emphasizes that autism is heterogeneous and cannot be reduced to one biological test or one predictable presentation.

My classroom experience made that limitation concrete. A diagnostic label cannot tell me what a particular learner understands, fears, enjoys, prefers, or needs in a specific moment. It cannot tell me whether the barrier is sensory, linguistic, technological, social, emotional, or simply a confusing set of directions. It cannot tell me what will change when the environment changes. Diagnosis is one source of evidence, but it is not the whole story.

This is also why labels can become dangerous when adults let them replace curiosity. They influence what people expect, which behaviors are interpreted as meaningful, and which students are assumed capable before they have a chance to show it. A diagnosis may create access to funding, services, and community while also inviting stigma or lowered expectations. The challenge is not to reject diagnostic knowledge. The challenge is to prevent it from becoming the only knowledge adults treat as credible.

\newpage
\reviewsection{Neurodiversity Became Real to Me Through Both Identity and Teaching}

\citet{ortega2013cerebralizing} connected with something I have felt for years: labels are not only clinical. They can become part of identity, community, resistance, and self-understanding. As someone who received an ADHD diagnosis later in life, I know the relief of finding language that explains patterns previously interpreted as laziness, inconsistency, or failure. A name can reduce shame. It can help a person reinterpret the past and ask for support without believing they are morally defective.

But I also understand Ortega's warning. Even an affirming brain-based identity can become reductive if it suggests that a person is only a brain or that everyone with the same label shares one political, cultural, or personal experience. That tension appears in classrooms too. A teacher can move away from deficit language and still create a new stereotype: the autistic child who is always visual, the student with ADHD who is always energetic, or the neurodivergent learner who is automatically gifted in one narrow way.

My students keep reminding me that difference does not arrive in neat categories. In one activity, a student needed a familiar laptop nearby while also building with blocks. The device could have been interpreted as distraction or refusal. Instead, I treated it as part of the student's regulation and allowed both forms of participation to coexist. The student remained connected to the activity. That moment taught me more than a checklist could: support does not always look like removing something. Sometimes it means allowing a learner to keep what helps them feel safe enough to participate.

For me, neurodiversity is most useful when it expands whose knowledge counts and protects variation within the spectrum. It should resist the idea that one form of communication, behavior, attention, or participation defines competence. It should make us more curious, not more certain.

\reviewsection{What the Media Helped Me Recognize About Representation}

Watching \textit{Autism: Insight from Inside} reminded me of how strongly public understanding depends on stories \citep{infobase2016inside}. I have experienced the power of being seen through the wrong story: the unreliable student, the overly emotional person, or the one who should simply try harder. I also know the relief of encountering a story that makes room for complexity rather than forcing a person into a deficit narrative. The film made that tension visible through multiple perspectives. Participants described how diagnosis can unlock services, language, funding, and community while also bringing stigma and lowered expectations. They showed that people with the same diagnosis can have very different sensory experiences, communication profiles, interests, and support needs. The film's emphasis on developing strengths shifted attention away from fixing a child and toward discovering possibility.

At the same time, I do not want one successful autistic person's story to become a new standard that everyone else must meet. Representation becomes harmful when one person is used to prove what an entire group is supposed to look like. The value of the film was not that it offered one model of autism. Its value was that it made room for different ways of experiencing the world and insisted that understanding begins by listening from the inside. That matters in teaching because students are also surrounded by stories. They learn which bodies are considered calm, which language is considered intelligent, which behavior is considered respectful, and which kinds of help are treated as weakness. Media can widen public understanding when it centers many voices. It can narrow understanding when one story becomes the measuring stick for everyone.

\reviewsection{What Students Taught Me About Behavior, Communication, and Design}

The National Education Association guide became most useful when I connected it to moments from my classroom \citep{nea2020guide}. During technology, coding, and robotics activities, I saw how movement, repetition, avoidance, silence, leaving a group, or refusing a material could be interpreted as defiance when the actual issue was overload, uncertainty, lack of control, or inaccessible directions. I also saw how quickly behavior changed when the environment changed.

When I placed only the needed materials on the table, modeled the first step, reduced the number of verbal directions, clarified peer roles, or offered a quieter starting point, students often re-entered the activity without a lecture, consequence, or reward. That is why the guide's framing of inclusion matters to me. Inclusion is not simply physical placement. Students must be valued members of the classroom, able to participate in instruction and relationships, and provided the supports they need.

The strongest question I can ask is not, \textit{``How do I stop this behavior?''} It is, \textit{``What might this student be communicating, and what barrier is present?''} That change in language changes the intervention. It moves me away from controlling the student and toward redesigning the conditions.

This also changed how I understand independence. Support should not mean doing the task for a student or making every decision on the student's behalf. The strongest scaffolds make the goal visible and the entry point accessible while preserving opportunities for choice, productive struggle, communication, revision, and repair. A visual sequence, partner role, sensory option, communication board, model, or quiet reset should reveal competence rather than replace it.

Communication and positive behavioral support therefore belong together. Both should increase access and agency, not enforce conformity.

\reviewsection{Why Credibility, Culture, and Identity Cannot Be Separated}

My lived experience also makes it impossible for me to think about developmental difference as an isolated variable. I am Latino, bilingual, neurodivergent, chronically ill, and shaped by trauma. Different institutions have interpreted the same traits differently depending on which part of me they noticed first. Sometimes persistence was seen as intensity. Sometimes difficulty with an inaccessible process was treated as lack of responsibility. Sometimes communication style mattered more than the content of what I was saying.

Students also move through the world with multiple identities. An autistic learner may be multilingual, racialized, gender-diverse, economically marginalized, physically disabled, or part of a family and culture with its own language for disability and support. Those identities affect who is noticed, who is believed, who receives an evaluation, which behavior is punished, and what services are available.

This is why lived experience cannot be treated as an optional add-on after the ``real'' academic analysis. Professional knowledge is incomplete without it. Family knowledge becomes incomplete only when institutions decide not to listen. Student behavior is incomplete evidence when adults refuse to consider context. A culturally responsive understanding requires student voice, family knowledge, attention to language and community, and awareness that the same behavior may be interpreted differently depending on who displays it.

\newpage
\reviewsection{Designing Better Evidence Through the Puzzle Plan and EQUITAS}

The readings, media, teaching experiences, and my own history come together in the Puzzle Plan. Diagnostic systems, school systems, and technology systems all create evidence about people. A shallow system records the label, error, delay, referral, absence, or behavior. A richer system also records task demands, sensory conditions, communication options, language, tool access, peer dynamics, prior knowledge, student preference, family knowledge, and what changed when support was offered.

I have seen the difference in real time. When a student struggled with a technology task, the shallow evidence was that the student did not complete it. The richer evidence included a confusing login sequence, too much spoken information, peer competition for the device, noise, unclear visual expectations, and insufficient processing time. Once those conditions changed, the student's participation changed too.

The quality of the evidence shapes the quality of the response. If the evidence is deficit-only, the intervention will usually ask the student to become easier for the system. If the evidence includes context and agency, educators can redesign the system to become more responsive to the student.

Through the EQUITAS lens, I now ask whether classroom structures provide equitable and accountable pathways for students to access learning, communicate needs, demonstrate competence, and exercise agency without erasing difference in order to belong. That question is not separate from the readings. It is what the readings become when they meet a real student, a real classroom, and a real life.

\reviewsection{Conclusion}

The world comes to understand developmental differences through labels, stories, relationships, institutions, environments, and the voices it chooses to treat as credible. My own life and teaching have shown me that a person can be described very differently depending on whether a system notices only the difficulty or also notices the conditions surrounding it.

A label can open access, but it becomes harmful when adults allow it to replace curiosity. In my future classroom, I want to combine diagnostic information with direct observation, student communication, family knowledge, sensory and environmental information, peer dynamics, lived experience, and collaboration with specialists. I want support to reveal competence rather than conceal agency.

The goal is not to make every learner communicate, move, focus, or participate in the same way. The goal is to create a classroom where students with autism and other complex support needs can communicate, make choices, demonstrate competence, belong, and be believed. Understanding developmental difference responsibly requires more than naming a condition. It requires staying curious about the person, collecting richer evidence, and changing the environment before deciding that the learner is the problem.

\newpage
\small
\bibliography{references}

\end{document}